%% ma: L12-B01-0048
%% lop: 12
%% chuong: 1
%% bai: 1
%% dang: 12.1.27
%% loai: TLN
%% muc: VD
%% dap_an: "33,3"
%% nguon: "Ảnh giáo viên gửi 05/11/2026 (M11): Kiểm tra giữa HK1 Toán 12 (2026-2027), THPT Hồ Thị Bi (TP.HCM), mã đề 1133, Phần III Câu 3"
%% kien_thuc: "Bài 1 (tính đơn điệu: hàm số giảm khi đạo hàm âm); Lớp 11 (đạo hàm của hàm đa thức, ý nghĩa cơ học của đạo hàm)"
%% trang_thai: cho-duyet
%% ly_do: "Dạng mới 12.1.27 đang đề xuất, chờ giáo viên duyệt"
%% ghi_chu: "Đề gốc không kèm đáp án; đã giải độc lập, đáp số 33,3. Đề gốc dẫn nguồn Bigalke et al., Mathematik Grundkurs ma-1 (Cornelsen 2016)."
\begin{ex}
Một tàu đổ bộ tiếp cận Mặt Trăng theo phương thẳng đứng và bắt đầu đốt cháy các tên lửa hãm ở độ cao $250\,\mathrm{km}$ so với bề mặt của Mặt Trăng. Giả sử trong khoảng $70$ giây đầu tiên kể từ khi đốt cháy các tên lửa hãm, độ cao $y$ của con tàu so với bề mặt của Mặt Trăng được mô hình hóa bởi hàm số
\[y=h(t)=-0{,}01t^3+1{,}1t^2-30t+250,\]
trong đó $t$ là thời gian tính bằng giây và $y$ là độ cao tính bằng kilômét.

Biết rằng vận tốc tức thời $v(t)=h'(t)$ của con tàu giảm dần trên khoảng thời gian $(t_1;t_2)$ với $0\le t_1<t_2\le70$. Khi đó, giá trị lớn nhất của hiệu số $t_2-t_1$ bằng bao nhiêu? (làm tròn kết quả đến hàng phần mười)
\shortans{33,3}
\loigiai{$v(t)=h'(t)=-0{,}03t^2+2{,}2t-30$, $v'(t)=-0{,}06t+2{,}2$.
$v$ giảm khi $v'(t)<0\Leftrightarrow t>\dfrac{110}3$.
Trên $[0;70]$, $v$ giảm trên khoảng $\left(\dfrac{110}3;70\right)$ nên $(t_2-t_1)_{\max}=70-\dfrac{110}3=\dfrac{100}3\approx33{,}3$.}
\end{ex}
